\subsection*{A single profinite seed for the adelic cubic Burau family}

\paragraph{BI.1 (Proved): the profinite exponent simplification.}
Let $e\in\widehat{\mathbb Z}=\prod_p\mathbb Z_p$ have
$e_3=1$ and $e_p=0$ for all $p\ne3$. This is the idempotent
of the three-adic direct factor. Define the actual profinite word
\[
 A_e=[[x^e,y^e],[y^e,z^e]],\qquad z=(xy)^{-1}.
\]
Then $A_e\in F''$, $e\equiv4\pmod{12}$, and
$q_{12}(A_e)=5I$ in the original signed representation.
The last assertion is the same finite calculation as for the
integer exponent $4$, since the peripheral images have orders
dividing $12$.

In the reduced cubic Burau representation over
$\mathcal O=\mathbb Z[\omega]$, $\omega^2+\omega+1=0$,
the matrices $A,B$ of FP.36 preserve the integral flag
\[
 \mathcal O(1,1,1)\subset
 \ker(1,\omega,\omega^2)\subset\mathcal O^3.
\]
In the unimodular basis
$v=(1,1,1)$, $w=(0,1,-\omega^2)$, $e_1=(1,0,0)$,
their matrices are
\[
 A=\begin{pmatrix}1&\omega&-\omega\\
                  0&-\omega&-\omega^2\\0&0&1\end{pmatrix},
 \qquad B=\operatorname{diag}(1,-\omega,1).
\]
Consequently
\[
 X_4=(AB)^3=I+(1-\omega)E_{13}.
\]
Its difference from $I$ is square-zero; the same holds for
$X_5$, which is its reversing-braid conjugate. Thus
$X_i^e=I$ for $i=4,5$ in every prime-to-three component
of $\mathcal O\otimes\widehat{\mathbb Z}$, simultaneously.
The corresponding three pair evaluations of $A_e$ are one,
as are all their peripheral $S_3$ conjugates: at least one
of the three powered peripheral variables is one, so one
of the two inner commutators is one.

The other two adjacent pair images are triangular in the
displayed flag or its reversing conjugate. Their derived
groups lie in the upper unitriangular group of degree three;
their second derived groups are therefore central and abelian.
At the prime above three, $\omega$ reduces to $1$ and
$A^2,B^2,C^2$ reduce to $I$. Every finite pure image there
is a three-group. These are exactly the nilpotent image
conditions used in the residual step of LI.9.

Finally $e$ is even. Under each peripheral conjugate of the
original orbifold quotient, one of $x,y,z$ has order two;
its $e$-th power is one. Hence $A_e$ belongs to the same
orbifold normal core required of the seed. The usual finite
intersection and compactness argument may therefore use
this single profinite word in place of varying integer
exponents $4M$ with $3\nmid M$. This simplifies the seed
choice only; it does not replace any of the lifting or
normal-core arguments in LI.9.
